\section{Introducing HQB}\label{sec:hqb}

We consider a construction which combines a $\KEM$ and a nominal group to
form a hybrid $\KEM$.

\subsection{Design}
%
$\HQB$ opts for a conservative construction while providing ease of
implementation given a common set of crypto tools: a KDF, an elliptic
curve group, and a post-quantum KEM.Relying directly on the SDH security
of the group allows us to inline the DH-KEM construction and avoid extra
layers of hashing, which is necessary elsewhere to achieve an IND-CCA KEM
from a group and then use that classical KEM as component of the hybrid
scheme. As group-based KEMs are less common in crypto libraries than
plain elliptic curve groups, this allows easier implementation and
deployment.

Including all ciphertexts of the component schemes lets us achieve
IND-CCA security~\cite{PKC:GiaHeuPoe18} in the pre- and post- quantum
settings without relying on C2PRI or other analysis of the PQ KEM
component. By including all the public key material of the component
schemes, we mitigate multi-target/multi-user
attacks~\cite{EPRINT:GlaHovSte25}, and generally cover our bases in terms
of `committing'~\cite{CCS:CreDaxMed24} to the public data pertaining to
the computation of the hybrid shared secret.
%

\begin{definition}[\hqb] \label{def:hqb}
    %
    Let $\nominalgroup = \nominalgroupexpanded$ be a nominal group, let
    $\kem$ be a key encapsulation mechanism and let $H$ be a hash function.
    %
    We define the \hqb \kem as depicted in \autoref{fig:hqb}.
    %
    \begin{gamefigure}{\hqb Framework.}{fig:hqb}
        \begin{games}{2}
            \begin{gameenv}{\algor $\keygen()$}
                \State $(\sk_1, \pk_1)  \sample \kem.\keygen()$
                \State $\sk_2 \sample \honestexponent$
                \State $\pk_2 \assign \nexp(\ngen, \sk_2)$
                \State $\pk \assign (\pk_1, \pk_2)$
                \State $\sk \assign (\sk_1, \sk_2, \pk)$
                \State \Return $(\sk, \pk)$
            \end{gameenv}
            \columnbreak
            \begin{gameenv}{\algor $\enc(\pk)$}
                \State $(\pk_1, \pk_2) \assign \pk$
                \State $k_1, c_1 \sample \kem.\enc(\pk_1)$
                \State $\sk_e \sample \honestexponent$
                \State $c_2 \assign \nexp(g, \sk_e)$
                \State $k_2 \assign \nexp(\pk_2, \sk_e)$
                \State $c \assign (c_1, c_2)$
                \State $k \assign H(k_1 \concat k_2 \concat c_1 \concat c_2 \concat \pk_1 \concat \pk_2 \concat \text{label})$
                \State \Return $(k, c)$
            \end{gameenv}
        \end{games}

        \begin{games}{2}
            \begin{gameenv}{\algor $\dec(c, \sk)$}
                \State $(\sk_1, \sk_2, \pk) \assign \sk$
                \State $(c_1, c_2) \assign c$
                \State $k_1 \assign \kem.\dec(c_1, \sk_1)$
                \State $k_2 \assign \nexp(c_2, \sk_2)$
                \State $\ifthen{k_1 \compare \bot}$
                \State \quad \Return $\bot$
                \State $k \assign H(k_1 \concat k_2 \concat c_1 \concat c_2 \concat \pk_1 \concat \pk_2 \concat \text{label})$
                \State \Return $k$
                \end{gameenv}
                \gamebreak
                \emptygame
            \end{games}
        \end{gamefigure}

    \end{definition}

    \subsection{Correctness}
    %
    The correctness of the HQB framework follows from the correctness of the
    \kem used to instantiate \hqb.
    %
    Formally, if \hqb is instantiated with a $\delta$-correct \kem, then the
    correctness of \hqb is defined as follows:
    %
    \[
        \Pr \left[\,  k \ne k' \, \left | \substack{(\sk, \pk) \sample \keygen(\,) \\ (k, c) \sample \enc(\pk) \\ k' \gets \dec(c, \sk) } \right. \right] \le \delta .
    \]
    %
